\begin{thebibliography}{99}
\bibitem{alexnet} A. Krizhevsky, I. Sutskever, and G. E. Hinton, ``\href{https://papers.nips.cc/paper_files/paper/2012/hash/c399862d3b9d6b76c8436e924a68c45b-Abstract.html}{ImageNet classification with deep convolutional neural networks},'' \emph{Advances in Neural Information Processing Systems, vol. 25}, 2012.
\bibitem{clip} A. Radford et al., ``\href{https://proceedings.mlr.press/v139/radford21a.html}{Learning transferable visual models from natural language supervision},'' \emph{Proc. ICML}, 2021, pp. 8748--8763.
\bibitem{siglip2} M. Tschannen et al., ``\href{https://arxiv.org/abs/2502.14786}{SigLIP 2: Multilingual vision-language encoders with improved semantic understanding, localization, and dense features},'' \emph{arXiv:2502.14786}, 2025.
\bibitem{pe} D. Bolya et al., ``\href{https://arxiv.org/abs/2504.13181}{Perception Encoder: The best visual embeddings are not at the output of the network},'' \emph{arXiv:2504.13181}, 2025.
\bibitem{sift1999} D. G. Lowe, ``\href{https://www.cs.ubc.ca/~lowe/papers/iccv99.pdf}{Object recognition from local scale-invariant features},'' \emph{Proc. ICCV}, 1999, pp. 1150--1157.
\bibitem{sift} D. G. Lowe, ``\href{https://www.cs.ubc.ca/~lowe/papers/ijcv04.pdf}{Distinctive image features from scale-invariant keypoints},'' \emph{International Journal of Computer Vision, vol. 60}, 2004, pp. 91--110.
\bibitem{bovw} J. Sivic and A. Zisserman, ``\href{https://www.robots.ox.ac.uk/~vgg/publications/2003/Sivic03/}{Video Google: A text retrieval approach to object matching in videos},'' \emph{Proc. ICCV}, 2003, pp. 1470--1477.
\bibitem{ridge} A. E. Hoerl and R. W. Kennard, ``\href{https://doi.org/10.1080/00401706.1970.10488634}{Ridge regression: Biased estimation for nonorthogonal problems},'' \emph{Technometrics, vol. 12}, 1970, pp. 55--67.
\bibitem{adamw} I. Loshchilov and F. Hutter, ``\href{https://arxiv.org/abs/1711.05101}{Decoupled weight decay regularization},'' \emph{Proc. ICLR}, 2019.
\bibitem{dropout2012} G. E. Hinton et al., ``\href{https://arxiv.org/abs/1207.0580}{Improving neural networks by preventing co-adaptation of feature detectors},'' \emph{arXiv:1207.0580}, 2012.
\bibitem{dropout} N. Srivastava et al., ``\href{https://jmlr.org/papers/v15/srivastava14a.html}{Dropout: A simple way to prevent neural networks from overfitting},'' \emph{Journal of Machine Learning Research, vol. 15}, 2014, pp. 1929--1958.
\bibitem{adam} D. P. Kingma and J. Ba, ``\href{https://arxiv.org/abs/1412.6980}{Adam: A method for stochastic optimization},'' \emph{Proc. ICLR}, 2015.
\bibitem{sgdr} I. Loshchilov and F. Hutter, ``\href{https://arxiv.org/abs/1608.03983}{SGDR: Stochastic gradient descent with warm restarts},'' \emph{Proc. ICLR}, 2017.
\bibitem{xavier} X. Glorot and Y. Bengio, ``\href{https://proceedings.mlr.press/v9/glorot10a.html}{Understanding the difficulty of training deep feedforward neural networks},'' \emph{Proc. AISTATS}, 2010, pp. 249--256.
\bibitem{heinit} K. He et al., ``\href{https://openaccess.thecvf.com/content_iccv_2015/html/He_Delving_Deep_into_ICCV_2015_paper.html}{Delving deep into rectifiers: Surpassing human-level performance on ImageNet classification},'' \emph{Proc. ICCV}, 2015, pp. 1026--1034.
\bibitem{vgg} K. Simonyan and A. Zisserman, ``\href{https://arxiv.org/abs/1409.1556}{Very deep convolutional networks for large-scale image recognition},'' \emph{Proc. ICLR}, 2015.
\bibitem{inception} C. Szegedy et al., ``\href{https://openaccess.thecvf.com/content_cvpr_2015/html/Szegedy_Going_Deeper_With_2015_CVPR_paper.html}{Going deeper with convolutions},'' \emph{Proc. CVPR}, 2015, pp. 1--9.
\bibitem{resnet} K. He et al., ``\href{https://openaccess.thecvf.com/content_cvpr_2016/html/He_Deep_Residual_Learning_2016_CVPR_paper.html}{Deep residual learning for image recognition},'' \emph{Proc. CVPR}, 2016, pp. 770--778.
\bibitem{transformer} A. Vaswani et al., ``\href{https://papers.nips.cc/paper_files/paper/2017/hash/3f5ee243547dee91fbd053c1c4a845aa-Abstract.html}{Attention is all you need},'' \emph{Advances in Neural Information Processing Systems, vol. 30}, 2017.
\bibitem{krr} C. Saunders, A. Gammerman, and V. Vovk, ``\href{https://eprints.soton.ac.uk/258942/}{Ridge regression learning algorithm in dual variables},'' \emph{Proc. ICML}, 1998, pp. 515--521.
\bibitem{pls1966} H. Wold, ``Estimation of principal components and related models by iterative least squares,'' \emph{Multivariate Analysis, P. R. Krishnaiah, Ed. Academic Press}, 1966, pp. 391--420.
\bibitem{pls} S. Wold, A. Ruhe, H. Wold, and W. J. Dunn III, ``\href{https://doi.org/10.1137/0905052}{The collinearity problem in linear regression. The partial least squares (PLS) approach to generalized inverses},'' \emph{SIAM Journal on Scientific and Statistical Computing, vol. 5}, 1984, pp. 735--743.
\bibitem{pca} K. Pearson, ``\href{https://doi.org/10.1080/14786440109462720}{On lines and planes of closest fit to systems of points in space},'' \emph{Philosophical Magazine, vol. 2}, 1901, pp. 559--572.
\bibitem{stacking} L. Breiman, ``\href{https://doi.org/10.1007/BF00117832}{Stacked regressions},'' \emph{Machine Learning, vol. 24}, 1996, pp. 49--64.
\bibitem{occlusion} M. D. Zeiler and R. Fergus, ``\href{https://arxiv.org/abs/1311.2901}{Visualizing and understanding convolutional networks},'' \emph{Proc. ECCV}, 2014, pp. 818--833.
\end{thebibliography}
